\section{条件配色}
\key{thresholds} 接受两个严格递增的阈值，\key{threshold colors}
接受三个颜色。颜色选择基于原始输入值，不基于四舍五入后的标签，
也不基于截断后的条长。

\begin{example}{示例 42：低、中、高三段配色}
\gradbarssetup{
  thresholds={60,80},
  threshold colors={
    lightred,gold,lightgreen
  }
}
\gradbar{59}

\gradbar{60}

\gradbar{79}

\gradbar{80}

\gradbar{92}
\end{example}

\begin{tabularx}{\linewidth}{@{}lX@{}}
\toprule
\multicolumn{1}{c}{\textbf{条件}} & \multicolumn{1}{c}{\textbf{本例颜色}}\\
\midrule
$v<60$ & lightred\\
$60\leq v<80$ & gold\\
$v\geq80$ & lightgreen\\
\bottomrule
\end{tabularx}

\begin{example}{示例 43：与目标和圆角叠加}
\gradbar[
  width=30mm,
  thresholds={60,80},
  rounded=2pt,
  target=80,
  label=end
]{72}
\end{example}

\note{启用条件配色后，所选颜色采用纯色填充，并优先于主题或
\key{negative color}。默认三色为 lightred、gold、lightgreen。
用 \key{thresholds={}} 关闭。阈值代表什么、数值越高是否越好，
应在表格标题或注释中说明。}
